\subsection{Zero Divisors and Spectra of Algebraic Elements}

\begin{definition} \label{def:algebra-zero-divisors-and-spectrum} \hlt{Zero Divisors and Spectrum}\\
Let $\mathcal{A}$ be an $\mathbb{F}$-algebra and $0\neq a\in\mathcal{A}$.
\begin{enumerate}[label=(\roman*)]
\setlength{\itemsep}{0pt}
\item $a$ is a \hlt{left zero divisor} iff $\exists b\in\mathcal{A}\setminus\{0\}$ with $ab=0$.
\item $a$ is a \hlt{right zero divisor} iff $\exists c\in\mathcal{A}\setminus\{0\}$ with $ca=0$.
\item If $a$ is algebraic, \hlt{spectrum} is $\operatorname{spec}_{\mathcal{A}}(a)\defeq\{r\in\mathbb{F}:m_a(r)=0\}$; elements are algebraic \hlt{eigenvalues} of $a$.
\end{enumerate}
\end{definition}

\begin{theorem} \label{thm:algebraic-invertible-iff-not-zero-divisor} \hlt{Invertible if and only if Not a Zero Divisor}\\
Let $a\in\mathcal{A}$ be algebraic. Then $a\in\mathcal{A}^{\times}$ iff $a$ is neither a left nor a right zero divisor.
\begin{proof}
\begin{enumerate}[label=\roman*.]
\setlength{\itemsep}{0pt}
\item[$(\Rightarrow)$] If $a\in\mathcal{A}^{\times}$ and $ab=0$, then $b=a^{-1}ab=0$; if $ca=0$, then $c=caa^{-1}=0$.\\
Hence $a$ is neither a left nor a right zero divisor.
\item[$(\Leftarrow)$] Prove the contrapositive. If $a\notin\mathcal{A}^{\times}$, Theorem \ref{thm:algebraic-element-annihilator-and-invertibility} gives $m_a(x)=xp(x)$ with $p\neq0$ and $\deg p<\deg m_a$.\\
Then $0=m_a(a)=ap(a)=p(a)a$, while minimality gives $p(a)\neq0$. Thus $a$ is both left, right zero divisor.
\end{enumerate}
\hlt{Structural Commentary}: For algebraic elements, noninvertibility forces a polynomial kernel witness.\\
\hlt{Key Insight}: Factoring $m_a$ by $x$ manufactures the zero-divisor witness $p(a)$.
\end{proof}
\end{theorem}

\begin{corollary} \label{cor:algebraic-one-sided-inverse-two-sided} \hlt{One-Sided Inverses Become Two-Sided}\\
Let $a\in\mathcal{A}$ be algebraic. If $\exists b\in\mathcal{A}$ with $ba=1_{\mathcal{A}}$ or $ab=1_{\mathcal{A}}$, then $a\in\mathcal{A}^{\times}$.
\begin{proof}
Assume $ba=1_{\mathcal{A}}$. If $ac=0$, then $c=bac=0$; if $ax=0$, then $x=bax=0$. Hence $a$ is neither a right nor a left zero divisor, so Theorem \ref{thm:algebraic-invertible-iff-not-zero-divisor} gives $a\in\mathcal{A}^{\times}$. The case $ab=1_{\mathcal{A}}$ is symmetric.\\
\hlt{Structural Commentary}: Algebraicity upgrades one-sided cancellation to two-sided invertibility.\\
\hlt{Key Insight}: A one-sided inverse rules out both zero-divisor obstructions.
\end{proof}
\end{corollary}

\begin{theorem} \label{thm:minimal-polynomial-roots-by-algebra-invertibility} \hlt{Roots Detected by Invertibility}\\
Let $a\in\mathcal{A}$ be algebraic and $r\in\mathbb{F}$. Then $m_a(r)=0\Longleftrightarrow a-r1_{\mathcal{A}}\notin\mathcal{A}^{\times}$. In particular, $a\in\mathcal{A}^{\times}\Longleftrightarrow0\notin\operatorname{spec}_{\mathcal{A}}(a)$.
\begin{proof}
\begin{enumerate}[label=\roman*.]
\setlength{\itemsep}{0pt}
\item[$(\Rightarrow)$] If $m_a(r)=0$, write $m_a=(x-r)q$ with $\deg q<\deg m_a$.\\
Then $0=m_a(a)=(a-r1_{\mathcal{A}})q(a)=q(a)(a-r1_{\mathcal{A}})$ and $q(a)\neq0$ by minimality, so $a-r1_{\mathcal{A}}\notin\mathcal{A}^{\times}$.
\item[$(\Leftarrow)$] If $m_a(r)\neq0$, then $\gcd(x-r,m_a)=1$, so $\exists u,v\in\mathbb{F}[x]$ with $u(x)(x-r)+v(x)m_a(x)=1$.\\
Evaluating at $a$ gives $u(a)(a-r1_{\mathcal{A}})=1_{\mathcal{A}}=(a-r1_{\mathcal{A}})u(a)$, hence $a-r1_{\mathcal{A}}\in\mathcal{A}^{\times}$, contradiction.
\end{enumerate}
Taking $r=0$ gives $a\in\mathcal{A}^{\times}\Longleftrightarrow m_a(0)\neq0\Longleftrightarrow0\notin\operatorname{spec}_{\mathcal{A}}(a)$.\\
\hlt{Structural Commentary}: Roots of $m_a$ are precisely the nonunit scalar shifts of $a$.\\
\hlt{Key Insight}: Bézout turns $\gcd(x-r,m_a)=1$ into an explicit inverse.
\end{proof}
\end{theorem}

\begin{theorem} \label{thm:algebra-polynomial-spectral-mapping} \hlt{Polynomial Spectral Mapping for Algebraic Elements}\\
Let $\mathbb{F}$ be algebraically closed, $a\in\mathcal{A}$ algebraic, and $p\in\mathbb{F}[x]$.\\
Then $p(a)$ is algebraic and $\operatorname{spec}_{\mathcal{A}}(p(a))=p(\operatorname{spec}_{\mathcal{A}}(a))\defeq\{p(r):r\in\operatorname{spec}_{\mathcal{A}}(a)\}$.
\begin{proof}
Since $\mathbb{F}[p(a)]\subseteq\mathbb{F}[a]$ and $\dim_{\mathbb{F}}\mathbb{F}[a]<\infty$, $p(a)$ is algebraic. If $p$ is constant, the claim is immediate.\\
If $\lambda\in\operatorname{spec}_{\mathcal{A}}(p(a))$, then $p(a)-\lambda1_{\mathcal{A}}\notin\mathcal{A}^{\times}$. Factor $p(x)-\lambda=c\prod_{j=1}^{m}(x-r_j)^{e_j}$ over $\mathbb{F}$.\\
Since a product of units is a unit, $a-r_j1_{\mathcal{A}}\notin\mathcal{A}^{\times}$ for some $j$. Theorem \ref{thm:minimal-polynomial-roots-by-algebra-invertibility} gives $r_j\in\operatorname{spec}_{\mathcal{A}}(a)$ and $\lambda=p(r_j)$.\\
Conversely, if $r\in\operatorname{spec}_{\mathcal{A}}(a)$, then $p(x)-p(r)=(x-r)q(x)$, so $p(a)-p(r)1_{\mathcal{A}}=(a-r1_{\mathcal{A}})q(a)$ is noninvertible; otherwise $a-r1_{\mathcal{A}}$ has a one-sided inverse and Corollary \ref{cor:algebraic-one-sided-inverse-two-sided} makes it invertible. Thus $p(r)\in\operatorname{spec}_{\mathcal{A}}(p(a))$.\\
\hlt{Structural Commentary}: Polynomial evaluation carries spectra by applying the same polynomial to roots.\\
\hlt{Key Insight}: Factor $p(x)-\lambda$ and detect a nonunit scalar shift of $a$.
\end{proof}
\end{theorem}

\begin{proposition} \label{prop:algebra-spectra-of-ab-and-ba} \hlt{Spectra of $ab$ and $ba$}\\
If $\dim_{\mathbb{F}}\mathcal{A}<\infty$, $\mathbb{F}$ is algebraically closed, and $a,b\in\mathcal{A}$, then $\operatorname{spec}_{\mathcal{A}}(ab)=\operatorname{spec}_{\mathcal{A}}(ba)$.
\begin{proof}
All elements are algebraic by Theorem \ref{thm:finite-dimensional-algebras-algebraic-elements}. Let $r\in\mathbb{F}^{\times}$.\\
If $ba-r1_{\mathcal{A}}\in\mathcal{A}^{\times}$, set $u\defeq a(ba-r1_{\mathcal{A}})^{-1}b-1_{\mathcal{A}}$; then $(ab-r1_{\mathcal{A}})u=r1_{\mathcal{A}}$, so $ab-r1_{\mathcal{A}}$ has a right inverse, hence is invertible by Corollary \ref{cor:algebraic-one-sided-inverse-two-sided}. Thus $r\notin\operatorname{spec}_{\mathcal{A}}(ba)\Rightarrow r\notin\operatorname{spec}_{\mathcal{A}}(ab)$, and symmetry gives equality off $0$.\\
If $ba\in\mathcal{A}^{\times}$ with inverse $y$, then $b(ay)=1_{\mathcal{A}}$ and $(yb)a=1_{\mathcal{A}}$, so $b,a\in\mathcal{A}^{\times}$ by Corollary \ref{cor:algebraic-one-sided-inverse-two-sided}; hence $ab\in\mathcal{A}^{\times}$. The converse is symmetric, so $0$ belongs to both spectra or neither.\\
\hlt{Structural Commentary}: $\operatorname{spec}(ab)\setminus\{0\}=\operatorname{spec}(ba)\setminus\{0\}$, and algebraicity resolves the $0$-case.\\
\hlt{Key Insight}: The resolvent identity transfers invertibility between $ba-r1_{\mathcal{A}}$ and $ab-r1_{\mathcal{A}}$ for $r\neq0$.
\end{proof}
\end{proposition}
